\section{Limitations and ethics}
\label{sec:limites}

\subsection{Technical delegability is not deployment}

ADI measures what an agent can do with the technology of October 2026, given the goal and access.
It does not measure what firms do, what regulators allow or what workers accept.
A task at L3 may never be delegated because the agent is too expensive relative to local wages, because the systems it would need are not connected, or because nobody is willing to be accountable for its output.
The index is therefore an upper bound on delegation, and the comparison with observed usage in Section~\ref{sec:validacion} is a comparison of a ceiling with a floor.
Readers who need a forecast of employment effects should combine ADI with adoption, prices and institutions, as \citet{acemoglu2024simple} does for exposure; the index alone supports none.
The reference technology is also a moving target.
We fix it by date and version the index (\texttt{v2026.1}) so that later releases can be compared with this one rather than silently replacing it.

\subsection{O*NET describes work in the United States}

The task statements, and the importance and frequency weights, come from surveys of United States workers.
The same occupation elsewhere may contain different tasks, in different proportions, with different tools.
The crosswalk to ISCO-08 carries the United States task mix to every country; it does not adapt it.
The problem is sharpest in agriculture.
O*NET has only 193 task statements that map to ISCO major group 6, skilled agricultural, forestry and fishery workers, which employs a large share of workers in low-income countries, and all 193 fall in the human-validation sample.
Since the delegation divide (H1) depends on that group, we rerun the country regression with the ESCO skill bank in place of O*NET tasks for major group 6 and for all groups; the slope with ESCO is \cifra{robustez.esco.h1.pendiente}, against \cifra{h1.pendiente} with O*NET.
Informal and subsistence work, which household surveys capture better than labour-force statistics by occupation, is treated in the microdata breakdowns rather than in the headline.

\subsection{Language models as raters}

The raters are language models from a single developer, and the usage data in Section~\ref{sec:validacion} come from the same developer's Anthropic Economic Index.
Three risks follow.
The first is bias: a model may rate tasks in its own domain, such as software, as more delegable than tasks it knows less well, and may share the blind spots of its training data about work outside rich countries.
The human validation on 1,500 tasks bounds this bias at the level of agreement reported in Section~\ref{sec:metodo}, but three raters with university degrees are not a sample of the world's workers.
The second is circularity: if the same family of models both rates delegability and generates the usage data, agreement between the two could reflect a shared view of the world rather than validation.
We address it by testing convergence against indices built by other methods and other raters \citep{eloundou2024gpts, felten2021occupational, gmyrek2025generative}, by comparing with the Agentic Adoption Index, which uses no model rating \citep{lee2026delegates}, and by reporting the Colombian wage dynamics, which involve no model at all.
The third is instability: the command-line interface exposes no temperature setting, so the same prompt can return different levels on different days.
The test--retest statistics in Section~\ref{sec:metodo} quantify this, and the consensus rule across models dampens it; the raw outputs of every pass are released so that others can recompute.
Model-based rating of social facts has a short literature \citep{horton2023large, argyle2023out, korinek2023language}, and the known tendency of models to favour their own outputs \citep{zheng2023judging, panickssery2024llm} is a reason to read every number here as a model's judgement, not a measurement of the world.

\subsection{Preregistration}

The hypotheses, rubric, model rule, sample seeds and analysis plan were registered before the first model call, under embargo until January 2027.
The registration constrains what we can report as confirmatory; everything else is labelled exploratory.
Deviations from the plan, including the change in the rater set described in Section~\ref{sec:metodo}, are dated in \texttt{docs/DESVIACIONES.md} and listed in the deviations appendix.
A preregistration does not make a measurement right; it makes it checkable.

\subsection{Ethics}

An index of delegability describes work in the voice of the machine: which parts of a job it could take, and which it could not.
That framing has costs.
It can reduce a worker to the residue of tasks an agent cannot do, and it can read a barrier such as trust or care as a temporary obstacle rather than as the point of the job.
We have tried to keep the four barriers visible rather than fold them into a single score, so that a reader can see why a task stays with a human and not only that it does.
The divide we name is also not symmetrical in its meaning.
In high-income countries, high delegability describes a labour market about to reorganise; in low-income countries, low delegability describes a labour market whose work agents cannot reach, which is not the same as protection.
Who is delegated, who delegates, and who is left outside the scope of delegation are distributional questions, and the atlas is meant to put them on the table, not to settle them.
We use no personal data; all inputs are public statistics, public classifications or anonymised household surveys used under their licences.
The human raters are informed adults with university degrees who agreed to the protocol; their individual ratings are released only in aggregate. % [VERIFY] compensation and ethics approval of the raters
